\documentclass[11pt]{article}

\usepackage[margin=1in]{geometry}
\usepackage{amsmath, amssymb, amsthm}
\usepackage{booktabs}
\usepackage{listings}
\usepackage{xcolor}
\usepackage{hyperref}

\lstdefinestyle{python}{
  language=Python,
  basicstyle=\ttfamily\small,
  keywordstyle=\color{blue!70!black},
  stringstyle=\color{green!50!black},
  commentstyle=\color{gray}\itshape,
  showstringspaces=false,
  breaklines=true,
  frame=single,
  rulecolor=\color{black!20},
  numbers=left,
  numberstyle=\tiny\color{gray},
}

\title{$\sigma$-Modulated UCB Exploration for Bilateral Matching}
\author{}
\date{}

\begin{document}
\maketitle

\section{$\sigma$-Modulated UCB Exploration}
\label{sec:sigma-ucb}

\subsection{Limitation of Count-Based UCB1}

The standard UCB1 selection rule for candidate-level exploration in the matching
system is
\begin{equation}
\widetilde{M}_v \;=\; M_v \;+\; c \cdot \sqrt{\frac{\ln t}{n_v + 1}},
\label{eq:ucb1}
\end{equation}
where $M_v$ is the analytical MapScore of candidate $v$, $n_v$ is the number of
times $v$ has been selected up to round $t$, and $c > 0$ is the exploration
constant. While Equation~\eqref{eq:ucb1} guarantees sublinear regret
$\mathcal{R}(T) = \mathcal{O}(\sqrt{T \log T})$, it imposes a uniform warmup:
every unseen candidate receives an unbounded bonus when $n_v = 0$, forcing the
system to try \emph{every} candidate at least once before exploiting. With pool
size $|\mathcal{V}| = 20$ and round budget $T = 60$, this warmup consumes $33\%$
of the budget regardless of how informative the initial priors are.

This penalty contradicts a design goal of our system: \emph{newcomers}
(candidates with no platform history) deserve an upfront exploration bonus, but
\emph{established candidates with a brief absence from the current task} do not.
Count-based UCB1 cannot distinguish these two cases because $n_v$ measures
only per-task selection history.

\subsection{Proposed Rule}

We replace Equation~\eqref{eq:ucb1} with a $\sigma$-modulated variant that uses
the candidate's posterior uncertainty $\bar{\sigma}_v$ as an exploration
intensity multiplier:
\begin{equation}
\widetilde{M}_v \;=\; M_v \;+\; c \cdot \bar{\sigma}_v \cdot
\sqrt{\frac{\ln t}{n_v + 1}},
\label{eq:sigma-ucb}
\end{equation}
where
$\bar{\sigma}_v = \frac{1}{K}\sum_{k=1}^{K} \sigma_v^k$
is the average posterior standard deviation across $v$'s $K$ capability
dimensions. Because $\sigma_v^k$ shrinks monotonically as the platform
accumulates feedback on $v$, $\bar{\sigma}_v$ encodes \emph{platform-wide}
information about $v$, not session-specific selection history.

Table~\ref{tab:bonus-intuition} illustrates the resulting exploration intensity
for representative candidate profiles at $c = 1$ and $t = 5$.

\begin{table}[h]
\centering
\small
\begin{tabular}{lccl}
\toprule
Candidate type & $\bar{\sigma}_v$ & Bonus & Behavior \\
\midrule
Brand-new (meta-knowledge source)   & 0.60 & 0.76 & strong upfront bonus \\
New (implicit-knowledge source)     & 0.40 & 0.51 & meaningful exploration \\
Established (explicit, few collabs) & 0.15 & 0.19 & mild bonus \\
Veteran (many past collaborations)  & 0.05 & 0.06 & near pure greedy \\
\bottomrule
\end{tabular}
\caption{Exploration bonus under $\sigma$-modulated UCB at $n_v = 0$, $c = 1$,
$t = 5$. Newcomers receive an order of magnitude more bonus than veterans.}
\label{tab:bonus-intuition}
\end{table}

Three properties follow directly:

\begin{enumerate}
\item \textbf{Upfront semantics.} New users enter the pool with high
$\bar{\sigma}_v$ from the L1 representation layer; they receive exploration
bonus immediately in round 1, without requiring repeated non-selection.

\item \textbf{No forced warmup.} An established candidate with $n_v = 0$ in the
current task does not receive an unbounded bonus, because $\bar{\sigma}_v$ is
small. Veterans compete on $M$ alone.

\item \textbf{Self-decaying exploration.} Once a newcomer is selected and
feedback is observed, the Bayesian posterior update shrinks $\sigma_v^k$, which
in turn reduces $\bar{\sigma}_v$ and the future exploration bonus. The system
does not over-explore.
\end{enumerate}

\subsection{Optional Sublinear Variant}

For settings where newcomer prominence should be amplified beyond linear
proportion, we propose a square-root variant:
\begin{equation}
\widetilde{M}_v \;=\; M_v \;+\; c \cdot \sqrt{\bar{\sigma}_v} \cdot
\sqrt{\frac{\ln t}{n_v + 1}}.
\label{eq:sqrt-sigma-ucb}
\end{equation}
This widens the gap between newcomer and veteran bonuses (Table~\ref{tab:variants}).

\begin{table}[h]
\centering
\small
\begin{tabular}{lcc}
\toprule
$\bar{\sigma}_v$ & Linear bonus (Eq.~\ref{eq:sigma-ucb}) & Square-root bonus (Eq.~\ref{eq:sqrt-sigma-ucb}) \\
\midrule
0.60 & 0.76 & 0.98 \\
0.40 & 0.51 & 0.80 \\
0.15 & 0.19 & 0.49 \\
0.05 & 0.06 & 0.28 \\
\bottomrule
\end{tabular}
\caption{Linear vs.\ square-root $\sigma$-modulation, $c = 1$, $t = 5$, $n_v = 0$.}
\label{tab:variants}
\end{table}

\subsection{Implementation}

The change requires modifying a single line of the existing candidate selector:

\begin{lstlisting}[style=python]
# UCB1 (original)
bonus = self.candidate_ucb_c * math.sqrt(math.log(t) / (n_v + 1))

# sigma-modulated UCB
avg_sigma_v = mean(cap.sigma for cap in candidate.capabilities)
bonus = self.candidate_ucb_c * avg_sigma_v \
        * math.sqrt(math.log(t) / (n_v + 1))
\end{lstlisting}

\section{Tiered Candidate Pool Design}
\label{sec:tiered-pool}

\subsection{Motivation: Coupling $\sigma$ and $\mu$ Noise}

A test pool with uniform $\sigma_v = 0.2$ and independently sampled $\mu$ noise
fails to exercise the $\sigma$-modulated UCB mechanism. To distinguish
$\sigma$-modulated UCB from count-based UCB1, the dataset must contain
candidates whose \emph{both} $\mu$-uncertainty and $\sigma$-uncertainty vary
together---reflecting the real-world fact that a platform with little history
on a user has both a noisy point estimate $\mu$ and low confidence $\sigma$.

We derive $\sigma$ from a combination of the L1 knowledge source and prior
collaboration history. Recall Proposition~C of the proposal, which states that
the posterior precision after $n$ updates satisfies
\begin{equation}
\sigma_v(n) \;\leq\; \frac{\sigma_{\text{obs}}^{\text{source}}}{\sqrt{n + 1}},
\end{equation}
where $\sigma_{\text{obs}}^{\text{source}} \in \{0.2, 0.4, 0.6\}$ for the
explicit, implicit, and meta-knowledge sources respectively.

\subsection{Tier Specification}

We define four candidate tiers, each parameterized by an initial $\sigma$ and a
$\mu$-noise standard deviation (Table~\ref{tab:tiers}).

\begin{table}[h]
\centering
\small
\begin{tabular}{llccc}
\toprule
Tier & Interpretation & $\sigma_v$ & $\mu$-noise & Pool fraction \\
\midrule
\texttt{veteran}      & Many past collaborations, well-known   & $\sim 0.05$ & $\mathcal{N}(0, 0.02)$ & 25\% \\
\texttt{mid}          & 2--3 past collaborations               & $\sim 0.15$ & $\mathcal{N}(0, 0.10)$ & 25\% \\
\texttt{new\_explicit} & Recently joined, detailed self-report & $\sim 0.30$ & $\mathcal{N}(0, 0.20)$ & 25\% \\
\texttt{new\_implicit} & Recently joined, inferred from chat   & $\sim 0.50$ & $\mathcal{N}(0, 0.30)$ & 25\% \\
\bottomrule
\end{tabular}
\caption{Four-tier candidate pool design with coupled $\sigma$ and $\mu$ noise.}
\label{tab:tiers}
\end{table}

For each candidate $v$ assigned to tier $\tau$, we draw
\begin{align}
\sigma_v &\sim \mathcal{U}(\sigma_\tau - \delta_\tau, \;\sigma_\tau + \delta_\tau), \\
\mu_v^{\text{init},k} &= \mathrm{clip}\bigl(\mu_v^{\text{true},k} + \varepsilon_k,\;0,\;1\bigr),
\quad \varepsilon_k \sim \mathcal{N}\bigl(0, \,\eta_\tau^2\bigr),
\end{align}
where $\sigma_\tau, \eta_\tau, \delta_\tau$ are the tier parameters from
Table~\ref{tab:tiers} and a small jitter $\delta_\tau$ (e.g., $0.02$--$0.08$)
prevents identical $\sigma$ values within a tier.

\subsection{Expected Algorithmic Behavior per Tier}

Under $\sigma$-modulated UCB with $c = 1$, $t = 5$:

\begin{itemize}
\item \textbf{Veterans} ($\sigma = 0.05$): exploration bonus $\approx 0.06$;
the system ranks them on near-true $\mu$ since their $\mu$-noise is also small.
\item \textbf{Mid-experience} ($\sigma = 0.15$): bonus $\approx 0.19$; mild
exploration if base $M$ is competitive.
\item \textbf{New explicit} ($\sigma = 0.30$): bonus $\approx 0.38$; $6\times$
a veteran's bonus, ensuring meaningful exploration despite noisy $\mu$.
\item \textbf{New implicit} ($\sigma = 0.50$): bonus $\approx 0.63$; $10\times$
a veteran's bonus. This is the tier most likely to harbor underestimated
high-ability candidates ("hidden gems") and the strongest test of the
anti-Matthew-effect property.
\end{itemize}


\end{document}
